% ============================================================
%  main.tex — IEEE Journal Paper
%  Zero-Shot Consumer Financial Complaint Classification
%  Using Open-Source LLMs on an HPC Cluster
% ============================================================
\documentclass[journal]{IEEEtran}

% ---- Packages -----------------------------------------------
\usepackage[T1]{fontenc}
\usepackage[utf8]{inputenc}
\usepackage{cite}
\usepackage{amsmath,amssymb}
\usepackage{graphicx}
\usepackage{booktabs}
\usepackage{array}
\usepackage{tabularx}
\usepackage{url}
\usepackage{hyperref}
\usepackage{multirow}
\usepackage{xcolor}
% balance handled by IEEEtriggeratref below
\usepackage{microtype}

% Tighten float spacing
\setlength{\floatsep}{8pt plus 2pt minus 2pt}
\setlength{\textfloatsep}{10pt plus 2pt minus 4pt}
\setlength{\intextsep}{8pt plus 2pt minus 2pt}
\setlength{\abovecaptionskip}{4pt}
\setlength{\belowcaptionskip}{0pt}

% Hyperref setup (IEEE-compatible)
\hypersetup{
  colorlinks=true,
  linkcolor=black,
  citecolor=black,
  urlcolor=black
}

% IEEE publication IDs are normally inserted after acceptance.
% \IEEEpubid{~\copyright~2026 IEEE}

% =============================================================
\begin{document}

% ---- Title --------------------------------------------------
\title{From Paid APIs to Local GPUs: Zero-Shot Consumer Financial\\
Complaint Classification with Open-Source LLMs}

% ---- Authors ------------------------------------------------
\author{Muhammet~Tan~and~Halit~Bak{\i}r%
\thanks{Manuscript received June 2026.}%
\thanks{M.~Tan and H.~Bak{\i}r are with the Department of Computer
Engineering, Sivas University of Science and Technology, Sivas, 58010,
Turkey (e-mail: muhammettan28@gmail.com; halit.bakir@sivas.edu.tr).}%
\thanks{Corresponding author: Halit Bak{\i}r.}%
}

\maketitle

% =============================================================
\begin{abstract}
Automated classification of consumer financial complaints has significant
implications for regulatory oversight and customer service prioritisation.
While recent studies have demonstrated strong zero-shot performance using
large proprietary language model (LLM) APIs, those approaches incur
substantial and recurring monetary costs that render them impractical for
resource-constrained researchers and institutions. In this paper we
investigate whether open-source, locally deployed LLMs can match or
approximate the classification quality of their paid-API counterparts,
entirely at zero API cost. We benchmark eight open-weight models
(four standard and four reasoning) against the Consumer Financial
Protection Bureau (CFPB) dataset, replicating an established five-category
zero-shot evaluation protocol on the same dataset. Models are deployed
via Ollama on an NVIDIA Tesla V100 GPU within the TRUBA ARF
high-performance computing (HPC) cluster, using SLURM for job
orchestration. We conduct two experiments: a balanced experiment
(200 samples per class, $N{=}1{,}000$) designed for protocol-level
comparison with the reference paper, and an imbalanced experiment
($N{=}1{,}000$, preprocessed natural class distribution) to assess
real-world robustness. Our best model,
\texttt{gpt-oss:20b}, achieves a weighted F1 of 0.7451 on the balanced
set, exceeding the reported paid GPT-4o score (F1~0.7358) under the
matched protocol with zero API spend.
Reasoning models form the strongest group on average, although
\texttt{phi4:14b} is a competitive low-latency standard-model exception. We further
document a complete, reproducible HPC inference pipeline and analyse
off-label prediction behaviour. Our results demonstrate that competitive
zero-shot complaint classification is achievable with open-source
infrastructure alone.
\end{abstract}

% ---- IEEE Keywords ------------------------------------------
\begin{IEEEkeywords}
consumer complaint classification, large language models, zero-shot
learning, open-source LLMs, high-performance computing, CFPB,
reasoning models, Ollama, SLURM, natural language processing
\end{IEEEkeywords}

% =============================================================
\section{Introduction}
\label{sec:intro}

The Consumer Financial Protection Bureau (CFPB) publishes a public
database of consumer complaints submitted against financial institutions
in the United States. The dataset spans dozens of product categories
and hundreds of thousands of narratives, making it one of the most
accessible and practically relevant corpora for natural language
processing (NLP) research in the financial domain. Accurately routing
and categorising these complaints has direct implications for
regulatory efficiency, bank compliance teams, and consumer advocacy
organisations: a miscategorised complaint may be forwarded to the wrong
department, delayed, or deprioritised
\cite{das2024explainable,liu2024framework,jondhale2024cfpb}.

Unlike many benchmark text-classification datasets, consumer complaint
narratives are not merely short topic labels embedded in clean prose.
They often contain informal descriptions of stressful financial events,
institution names, dates, partial account details, and overlapping
product issues. A single complaint may mention credit reporting, debt
collection, and a credit card dispute in the same narrative, while the
supervised label assigns only one product category. This ambiguity makes
the task more realistic than many standard NLP benchmarks and also more
valuable as a test bed for instruction-following models: a useful system
must map noisy, human-written narratives to a constrained regulatory
taxonomy without inventing new labels.

The practical motivation is equally strong. Complaint triage is a
high-volume activity in which automation can reduce manual screening
time, support earlier detection of recurring institutional problems,
and provide consistent first-pass routing. At the same time, financial
complaint text is sensitive. Even when datasets are public, a deployable
classification workflow for banks, regulators, or research laboratories
may need to operate without sending narratives to third-party services.
This makes local inference, auditability, and reproducible infrastructure
more than implementation details; they are part of the methodological
problem.

Traditional approaches to complaint classification relied on
bag-of-words representations, support vector machines, and logistic
regression. The emergence of BERT \cite{devlin2019bert} and its
domain-adapted successors marked a step change, enabling much richer
semantic representations and pushing classification accuracy above
90\,\% in fine-tuned settings \cite{vairetti2024bert}. More recently,
the advent of instruction-following LLMs such as GPT-3 \cite{brown2020gpt3}
and its successors opened the possibility of \emph{zero-shot}
classification: a model is simply given the candidate categories in
a prompt and is asked to assign the most appropriate one, requiring
no task-specific training data at all
\cite{kojima2022zeroshotcot,madanchian2024genai}.

For consumer complaint analysis, this zero-shot setting is attractive
because financial taxonomies evolve and labelled data may be costly to
curate for every new jurisdiction, product definition, or institutional
workflow. However, zero-shot classification also shifts the main
engineering burden from model training to prompt adherence, output
parsing, and deployment cost. A model that returns a plausible but
unsupported category such as \texttt{credit\_collection} may look
linguistically competent while still failing the regulatory routing
task. Therefore, evaluating off-label behaviour is essential alongside
accuracy and F1.

Roumeliotis~et~al.\ \cite{roumeliotis2025think} conducted a
comprehensive zero-shot benchmark of 14 state-of-the-art proprietary
and commercial models — including GPT-4o, GPT-4.5, Claude~3.5/3.7,
Gemini~2.0, o1, o3-mini, and DeepSeek (API) — on the same CFPB
dataset and five-category task. They reported that the best model,
OpenAI's o1, achieved an accuracy of 0.777 and an F1 of 0.7764, at a
cost of \$24.24 per thousand predictions. Their study provided strong
evidence that reasoning-enabled LLMs can outperform standard chat
models on this structured classification task. However, a critical
limitation of their work is that \emph{every} model was accessed
through a paid commercial API, creating a barrier to reproduction
and deployment in low-budget or data-sensitive environments.

This leaves an important open question: can open-weight models, served
locally on shared research infrastructure, deliver competitive
zero-shot complaint classification without API access? The answer is
not obvious. Smaller local models are cheaper to run but may be weaker
at following strict JSON-style instructions. Reasoning models may
improve disambiguation and output discipline, but they also generate
longer traces and increase latency. In addition, prior balanced
benchmarks do not show whether the same ranking survives under the
natural class imbalance of the CFPB corpus, where a majority
\texttt{credit\_reporting} class can inflate aggregate metrics.

The present paper addresses this gap directly. Our contributions are:

\begin{enumerate}
  \item We benchmark eight open-source, open-weight LLMs —
        \texttt{gemma2:9b}, \texttt{mistral:7b}, \texttt{llama3.1:8b},
        \texttt{phi4:14b}, \texttt{deepseek-r1:8b},
        \texttt{deepseek-r1:14b}, \texttt{qwen3:8b}, and
        \texttt{gpt-oss:20b} — on the same five-category CFPB
        zero-shot task, at \textbf{zero API cost}.

  \item We document a complete, reproducible inference pipeline on the
        TRUBA ARF HPC cluster using Ollama and SLURM, including
        all deployment parameters needed to replicate our results.

  \item We conduct a second experiment using the natural (imbalanced)
        class distribution, evaluating whether performance patterns
        observed under artificial balance hold in a more realistic
        scenario.

  \item We provide a direct quantitative comparison with the paid-API
        results of Roumeliotis~et~al.\ \cite{roumeliotis2025think},
        showing that our best open-source model exceeds the reported
        GPT-4o score under the same evaluation protocol and remains
        competitive with more expensive API-based alternatives.

  \item We analyse off-label prediction behaviour — instances where
        models invent labels outside the five prescribed categories —
        as a proxy for instruction-following reliability.
\end{enumerate}

The remainder of this paper is structured as follows.
Section~\ref{sec:related} surveys the related literature.
Section~\ref{sec:data} describes the dataset and preprocessing
pipeline. Section~\ref{sec:setup} details the experimental setup.
Section~\ref{sec:results} presents results for both experiments.
Section~\ref{sec:discussion} discusses the findings.
Section~\ref{sec:conclusion} concludes with directions for future work.


% =============================================================
\section{Related Work}
\label{sec:related}

\subsection{From Rule-Based Systems to Transformer Models}

The challenge of automatically classifying consumer complaints has a
rich history that mirrors the broader trajectory of NLP research. Early
systems relied on keyword matching and hand-crafted feature engineering
paired with classical classifiers such as support vector machines and
na{\"i}ve Bayes. While effective for narrow, well-defined complaint
taxonomies, these approaches proved brittle in the face of the
linguistic diversity found in real-world financial complaint corpora
\cite{raju2022bertopic,jondhale2024cfpb}.

The introduction of BERT \cite{devlin2019bert} triggered a significant
paradigm shift. Contextual word embeddings, obtained by pre-training on
large general-domain corpora and then fine-tuning on downstream tasks,
dramatically reduced the gap between human and machine performance on
text classification. Raju~et~al.\ \cite{raju2022bertopic} applied
BERTopic \cite{grootendorst2020bertopic}, FinBERT, DistilBERT, and
RoBERTa to a corpus of 684,653 CFPB complaints, demonstrating that
transformer-based topic models could discover latent complaint themes
without manual annotation. Vairetti~et~al.\ \cite{vairetti2024bert}
used the Spanish BERT variant BETO for complaint prioritisation,
achieving 92.1\,\% accuracy, illustrating that the BERT paradigm
transfers well across languages and regulatory contexts. Das~et~al.\
\cite{das2024explainable} pushed further by coupling transformer-based
classifiers with explainability techniques, enabling practitioners to
understand \emph{why} a complaint was assigned a given category — a
requirement for regulatory compliance. Liu~et~al.\ \cite{liu2024framework}
proposed a structured LLM-based categorisation framework (CCF) with
real-time monitoring, demonstrating that modern LLMs can serve as
end-to-end complaint classification pipelines with minimal domain
customisation.

\subsection{The LLM Era and the Zero-Shot Paradigm}

The release of GPT-3 \cite{brown2020gpt3} introduced the few-shot and
zero-shot paradigms, in which a large pre-trained language model can
perform novel tasks from a natural language description alone. This
capability fundamentally changed the cost calculus of NLP system
development: instead of collecting thousands of labelled examples and
retraining a model, practitioners could now describe a task in a
prompt and obtain reasonable predictions immediately. Zhao~et~al.\
\cite{zhao2023survey} provide a comprehensive survey of the capabilities
and limitations of modern LLMs, noting that instruction-following
performance scales strongly with model size and training quality.

In the financial complaint domain, Correa~et~al.\ \cite{correa2024genai}
demonstrated that generative AI can simultaneously classify, summarise,
and draft responses to CFPB complaints, achieving 88\,\% classification
accuracy in a zero-shot setting. This result was notable because it
required no fine-tuning whatsoever, confirming that the rich language
representations encoded in large-scale pre-training are sufficient for
this task. Extending the zero-shot paradigm to service industries,
Roumeliotis~et~al.\ \cite{roumeliotis2024tourism} compared GPT-4 Omni
and BERT for customer review classification in tourism, finding that
GPT-4o achieved 67\,\% accuracy in a zero-shot setting versus 60.6\,\%
for fine-tuned BERT — an early indication that zero-shot LLMs can
challenge fine-tuned specialised models. Li~et~al.\
\cite{li2024customerservice} further showed, through a comprehensive
customer service benchmark (TeleEval-CS), that fine-tuned open-source
LLMs can match proprietary counterparts when evaluated on
domain-specific dialogue tasks, motivating our investigation of
open-weight models.

A practically important contextual finding is that the CFPB complaint
corpus itself is changing. Shin~et~al.\ \cite{shin2025llmwritten}
analysed over 1.1 million CFPB complaints and found that approximately
9.8\,\% of complaints submitted after the public release of ChatGPT
exhibit strong evidence of LLM assistance. Liang~et~al.\
\cite{liang2025cfpb} independently estimated that roughly 18\,\% of
recent complaints are LLM-assisted. These findings underscore the
increasing entanglement of generative AI and financial consumer
protection, making robust automated classification even more important.

\subsection{Reasoning Models and the State of the Art}

The most recent development in this space has been the introduction of
chain-of-thought reasoning models \cite{wei2022chainofthought}, which
generate intermediate reasoning steps before arriving at a final answer.
DeepSeek-AI \cite{deepseekllm2024} demonstrated that open-source models
trained with careful scaling laws can match GPT-3.5 tier performance,
establishing the viability of open-weight alternatives before the
reasoning model era. Building on this foundation, Guo~et~al.\
\cite{guo2025deepseekr1} introduced DeepSeek-R1, an open-weight model
trained via reinforcement learning to produce structured reasoning
traces, achieving performance competitive with proprietary frontier
models on a range of benchmarks. Similarly, the
Qwen3 series \cite{qwen2025qwen3} introduced an explicit thinking mode
that can be toggled per query.

Roumeliotis~et~al.\ \cite{roumeliotis2025think} brought this reasoning
paradigm directly to the complaint classification setting. Testing
14 models on a balanced sample of 1,000 CFPB complaints across five
categories, they found that o1 (0.777 accuracy, F1~0.7764) and
GPT-4.5 (0.774, F1~0.7763) outperformed their non-reasoning
counterparts, and that reasoning models as a class exhibited fewer
off-label predictions. Their study is the closest antecedent to the
present work and provides the benchmark against which our results
are measured.

Concurrently, Bulgakov and Turchin \cite{bulgakov2026deepseek} explored
DeepSeek-R1-Distill-LLaMA-8B for binary prediction of monetary relief
outcomes from CFPB complaints in a zero-shot setting, demonstrating
that distilled reasoning models retain competitive performance while
fitting on a single consumer-grade GPU. G{\"u}ll{\"u}~et~al.\
\cite{gullu2025llama} fine-tuned LLaMA-3.1-8B on the same five-category
CFPB task used in our study, reaching 89.99\,\% accuracy and
establishing an approximate fine-tuning performance ceiling for that
model family.

\subsection{The Open-Source Gap}

Despite these advances, a critical gap remains in the literature. Every
paper benchmarking LLMs on CFPB complaint classification has relied
either on paid commercial APIs or on fine-tuned models that require
substantial labelled training data. Neither approach is accessible to
practitioners operating under strict data-governance constraints,
limited budgets, or air-gapped computing environments. The present
paper fills this gap by restricting itself entirely to open-weight
models deployed locally on shared HPC infrastructure, thereby
demonstrating that competitive zero-shot complaint classification is
achievable without any API cost or proprietary model access.


% =============================================================
\section{Dataset and Preprocessing}
\label{sec:data}

\subsection{Data Source}

We use the CFPB Consumer Complaints dataset originally published by
Tiwari \cite{tiwari2021cfpb} on Kaggle, which is a processed subset
of the public CFPB Consumer Complaint Database. The dataset covers
complaints submitted between March~2020 and March~2021 and contains
three columns: an unnamed index, \texttt{product} (the category label),
and \texttt{narrative} (the free-text complaint). The raw file
comprises 162,421 records consolidated to five target product categories:
\texttt{retail\_banking}, \texttt{credit\_reporting},
\texttt{credit\_card}, \texttt{mortgages\_and\_loans}, and
\texttt{debt\_collection}. This is the same dataset and label set used
by Roumeliotis~et~al.\ \cite{roumeliotis2025think}, ensuring
direct comparability.

\subsection{Class Distribution}

The original data are heavily imbalanced. After removing records with
empty \texttt{product} or \texttt{narrative} fields (10 rows removed,
162,411 remaining) and discarding narratives exceeding 500 characters
(resulting in 103,915 usable records), the class distribution shifts
to: \texttt{credit\_reporting}~63.6\,\%,
\texttt{debt\_collection}~14.1\,\%,
\texttt{mortgages\_and\_loans}~7.7\,\%,
\texttt{credit\_card}~7.5\,\%, and
\texttt{retail\_banking}~7.1\,\%. Before the length filter, the same
five-category corpus contains 162,411 records with
\texttt{credit\_reporting} accounting for 56.1\,\% of the data.

\subsection{Preprocessing Pipeline}

The preprocessing steps, implemented in Python using pandas, are applied
in the following order to every record:

\begin{enumerate}
  \item Drop rows with empty \texttt{product} or \texttt{narrative}.
  \item Retain only records belonging to the five target categories.
  \item Remove narratives longer than 500 characters, reducing processing
        time and GPU memory pressure while avoiding truncation artifacts
        during inference.
  \item Standardise narrative text: convert to lowercase, remove special
        characters, and normalise whitespace.
\end{enumerate}

\subsection{Experimental Subsets}

Two experimental subsets are drawn from the preprocessed corpus.

\textbf{Experiment~A (Balanced):} 200 records are sampled uniformly
at random from each of the five categories using stratified
undersampling with a fixed random seed of 42, yielding
$N = 1{,}000$ records. This follows the evaluation protocol of
Roumeliotis~et~al.\ \cite{roumeliotis2025think}; because the original
sample identifiers are not publicly released, the comparison should be
interpreted as protocol-matched rather than instance-identical.

\textbf{Experiment~B (Imbalanced):} 1,000 records are sampled using
stratified sampling that preserves the post-filter natural class
distribution:
\texttt{credit\_reporting}~636 (63.6\,\%),
\texttt{debt\_collection}~141 (14.1\,\%),
\texttt{mortgages\_and\_loans}~77 (7.7\,\%),
\texttt{credit\_card}~75 (7.5\,\%), and
\texttt{retail\_banking}~71 (7.1\,\%).
This second experiment, absent from the reference paper, assesses
how class imbalance affects both accuracy estimates and per-class
performance patterns.


% =============================================================
\section{Experimental Setup}
\label{sec:setup}

\subsection{Infrastructure}

All inference experiments are conducted on the TRUBA ARF cluster,
the Turkish national high-performance computing facility. Jobs are
submitted via the SLURM scheduler to the \texttt{akya-cuda} partition.
Each job is allocated a single NVIDIA Tesla V100-SXM2-16\,GB GPU
(16\,GB VRAM), 10~CPU cores, and a 2-hour wall-clock time limit.
Compute nodes have no outbound internet access; all models are pulled
in advance on the login node (\texttt{arf-ui1}) and stored in the
shared \texttt{/arf/home} filesystem, which is mounted read-only on
compute nodes. This offline-inference configuration mirrors real-world
deployments in air-gapped or regulated environments.

Model inference is handled by Ollama~v0.30.6, a self-contained
serving binary that manages model loading, KV-cache management,
and the REST API. The Ollama server is started within the SLURM
job, and the Python inference script communicates with it via the
local REST endpoint.

\subsection{Models}

Table~\ref{tab:models} lists all eight models. Four are standard
(non-reasoning) instruction-tuned models; four are reasoning models
that generate chain-of-thought traces before producing a final answer.
The evaluated Ollama builds use 4-bit quantisation (Q4\_K\_M GGUF
format, Ollama default), which reduces VRAM requirements significantly
while maintaining competitive quality.
The largest model, \texttt{gpt-oss:20b}, requires 13.0\,GB VRAM and
fits within a single V100.
Model names in Table~\ref{tab:models} are the exact Ollama tags used
in the experiments. The family column records the upstream lineage or
registry metadata associated with each tag and should not be interpreted
as equivalence to a proprietary API model with a similar provider name.
All artifacts were obtained through \texttt{ollama pull <tag>} with no
task-specific fine-tuning or local weight modification.
\texttt{gpt-oss:20b} is an open-weight reasoning model released by
OpenAI and distributed through the Ollama registry; it requires no API
key and its weights are publicly accessible via
\texttt{ollama pull gpt-oss:20b}.
The ``OpenAI'' entry in the Family column reflects the model's
originating developer, not access through any proprietary API or
commercial service.

\begin{table}[!t]
  \renewcommand{\arraystretch}{1.3}
  \caption{Open-Source Models Evaluated (Q4\_K\_M Quantisation Used in This Study)}
  \label{tab:models}
  \centering
  \footnotesize
  \setlength{\tabcolsep}{3pt}
  \begin{tabular}{@{}lllrr@{}}
    \toprule
    \textbf{Model} & \textbf{Family} & \textbf{Type} & \textbf{Params} & \textbf{VRAM (q4)} \\
    \midrule
    \texttt{gemma2:9b}        & Google     & Standard  & 9B   & 5.4\,GB \\
    \texttt{mistral:7b}       & Mistral    & Standard  & 7B   & 4.4\,GB \\
    \texttt{llama3.1:8b}      & Meta       & Standard  & 8B   & 4.9\,GB \\
    \texttt{phi4:14b}         & Microsoft  & Standard  & 14B  & 9.1\,GB \\
    \midrule
    \texttt{deepseek-r1:8b}   & DeepSeek   & Reasoning & 8B   & 5.2\,GB \\
    \texttt{deepseek-r1:14b}  & DeepSeek   & Reasoning & 14B  & 9.0\,GB \\
    \texttt{qwen3:8b}         & Alibaba    & Reasoning & 8B   & 5.2\,GB \\
    \texttt{gpt-oss:20b}      & OpenAI     & Reasoning & 20B  & 13.0\,GB \\
    \bottomrule
  \end{tabular}
\end{table}

\subsection{Prompt Design}

A single, model-agnostic prompt is used for all models and all
experiments. The prompt follows the five-label JSON-output template
used by Roumeliotis~et~al.\ \cite{roumeliotis2025think} to ensure
comparability:

\begin{quote}
\textit{You are an AI assistant specializing in consumer complaint
classification. Your task is to analyze a consumer complaint and
classify it into the most appropriate category from the predefined
list: ["retail\_banking", "credit\_reporting", "credit\_card",
"mortgages\_and\_loans", "debt\_collection"]. Provide your final
classification in the following JSON format without explanations:
\texttt{\{"product": "chosen\_category\_name"\}}. Complaint: [NARRATIVE]}
\end{quote}

The temperature parameter is set to 0 for all models to reduce
sampling variability and improve reproducibility. No explicit
\texttt{top\_p}, \texttt{max\_tokens}, stop-sequence, or model-specific
thinking-mode flags are set; Ollama defaults are used for all other
decoding parameters. Although the prompt asks for JSON without
explanations, some reasoning-oriented builds may still emit
intermediate text before the final answer. The inference script
therefore extracts the final JSON object from the response by searching
for the last valid \texttt{\{"product": "..."\}} occurrence using a
regular expression, accounting for models that prepend reasoning traces
before the answer.

The inference script was implemented using only Python standard-library
modules (\texttt{csv}, \texttt{json}, \texttt{urllib}, \texttt{time},
\texttt{argparse}), eliminating the need for any \texttt{pip}-installed
dependencies on the compute node. Each request is retried once on
HTTP failure before being recorded as an error; errors are included in
the final count and do not silently inflate accuracy. The raw model
response is retained alongside the parsed label in the output CSV to
allow post-hoc auditing of all parsing decisions. The
\texttt{gpt-oss:20b} model required a 4-hour SLURM wall-clock limit
due to its larger size; all other models completed inference within
the standard 2-hour limit.

\subsection{Evaluation Metrics}

Following the reference paper, we compute four primary metrics over
all $N{=}1{,}000$ predictions: \emph{exact-match accuracy},
\emph{weighted precision}, \emph{weighted recall}, and
\emph{weighted F1-score}. We additionally report \emph{macro-F1},
which weights all classes equally and is especially informative for
the imbalanced experiment. Weighted variants weight each class by its
sample size, making them appropriate for comparison with the reference
paper but potentially optimistic under class imbalance. We additionally
record the \emph{mean
latency} (seconds per prediction, wall-clock time) and the number
of \emph{off-label} predictions — responses where the model returned
a category name not in the prescribed five, or failed to return
valid JSON entirely. Off-label predictions are counted as errors.
The API cost for all experiments is \$0. Compute time is supplied by
an institutional HPC allocation rather than a commercial LLM API.


% =============================================================
\section{Results}
\label{sec:results}

\subsection{Experiment A: Balanced Evaluation}

Table~\ref{tab:results_balanced} presents the full results for
Experiment~A. The \texttt{gpt-oss:20b} reasoning model achieves the
highest accuracy (0.7450) and F1 (0.7451), followed by
\texttt{deepseek-r1:14b} (F1~0.7162) and the standard model
\texttt{phi4:14b} (F1~0.7095). Three reasoning models occupy the top four
positions (\texttt{gpt-oss:20b} rank~1, \texttt{deepseek-r1:14b}
rank~2, \texttt{qwen3:8b} rank~4), with the standard
\texttt{phi4:14b} ranked third. Notably,
\texttt{deepseek-r1:8b} (rank~6) falls below the standard
\texttt{gemma2:9b} (rank~5), indicating that model size interacts
with the reasoning benefit. The two weakest models are
\texttt{mistral:7b} (F1~0.5842) and \texttt{llama3.1:8b}
(F1~0.5958), both of which also produced the highest off-label
counts (54 and 57, respectively).

\begin{table}[!t]
  \renewcommand{\arraystretch}{1.3}
  \caption{Experiment A Results — Balanced ($N = 1{,}000$)}
  \label{tab:results_balanced}
  \centering
  \scriptsize
  \setlength{\tabcolsep}{2pt}
  \resizebox{\columnwidth}{!}{%
  \begin{tabular}{@{}llrrrrrrc@{}}
    \toprule
    \textbf{Model} & \textbf{Type} & \textbf{Acc.} & \textbf{Prec.} & \textbf{Rec.} & \textbf{F1} & \textbf{Macro-F1} & \textbf{Lat. (s)} & \textbf{OL} \\
    \midrule
    \texttt{gpt-oss:20b}      & R & 0.745 & 0.764 & 0.745 & 0.745 & 0.745 & 2.59 & 0 \\
    \texttt{deepseek-r1:14b}  & R & 0.716 & 0.733 & 0.716 & 0.716 & 0.716 & 5.92 & 0 \\
    \texttt{phi4:14b}         & S & 0.709 & 0.736 & 0.709 & 0.710 & 0.710 & 0.55 & 1 \\
    \texttt{qwen3:8b}         & R & 0.707 & 0.739 & 0.707 & 0.707 & 0.707 & 3.89 & 0 \\
    \texttt{gemma2:9b}        & S & 0.676 & 0.717 & 0.676 & 0.673 & 0.673 & 0.66 & 1 \\
    \texttt{deepseek-r1:8b}   & R & 0.675 & 0.709 & 0.675 & 0.674 & 0.674 & 4.64 & 1 \\
    \texttt{mistral:7b}       & S & 0.596 & 0.676 & 0.596 & 0.584 & 0.584 & 0.24 & 54 \\
    \texttt{llama3.1:8b}      & S & 0.587 & 0.697 & 0.587 & 0.596 & 0.596 & 0.57 & 57 \\
    \bottomrule
    \multicolumn{9}{l}{\footnotesize R=Reasoning, S=Standard, OL=Off-label, Lat.=Mean latency}
  \end{tabular}
  }
\end{table}

The \texttt{OL} column in Table~\ref{tab:results_balanced} is
particularly informative. The three zero-off-label models
(\texttt{gpt-oss:20b}, \texttt{deepseek-r1:14b}, and
\texttt{qwen3:8b}) are all reasoning models, whereas the two models
with severe format-following failures are standard models:
\texttt{llama3.1:8b} produces 57 off-label predictions and
\texttt{mistral:7b} produces 54. Aggregated by model type, the
reasoning group attains a mean F1 of 0.710 with only one off-label
prediction across 4,000 outputs, while the standard group attains a
mean F1 of 0.641 with 113 off-label predictions. This makes
instruction adherence, not only classification accuracy, a central
part of the performance gap. At the same time, \texttt{phi4:14b}
is an important exception among standard models: it reaches F1~0.7095
with only one off-label prediction and a mean latency of 0.55~s,
making it the strongest low-latency alternative.

Because all models are evaluated on the same 1,000 balanced examples,
we also perform post-hoc exact McNemar tests on paired correctness
outcomes for the leading model. The advantage of
\texttt{gpt-oss:20b} is statistically significant against
\texttt{deepseek-r1:14b} ($p=0.0106$), \texttt{phi4:14b}
($p=0.0030$), \texttt{qwen3:8b} ($p=0.0008$), and
\texttt{gemma2:9b} ($p<10^{-7}$). All four comparisons remain below
a Bonferroni-adjusted threshold of $\alpha=0.0125$. These tests
support the ranking observed in Table~\ref{tab:results_balanced},
while still leaving multi-seed resampling as an important direction
for future work.
For example, \texttt{gpt-oss:20b} is correct on 88 instances where
\texttt{phi4:14b} is incorrect, while the reverse occurs on 52
instances; the corresponding counts are 75 versus 46 for
\texttt{deepseek-r1:14b}, 81 versus 43 for \texttt{qwen3:8b}, and
115 versus 46 for \texttt{gemma2:9b}.
By contrast, adjacent high-performing models are not clearly
separated: \texttt{phi4:14b} is correct on 71 instances where
\texttt{qwen3:8b} is incorrect, while the reverse occurs on 69
instances ($p=0.9327$). The corresponding comparisons are also
non-significant for \texttt{deepseek-r1:14b} versus \texttt{phi4:14b}
(76 versus 69, $p=0.6184$) and \texttt{deepseek-r1:14b} versus
\texttt{qwen3:8b} (67 versus 58, $p=0.4744$), indicating that their
rank ordering should be interpreted cautiously.

For Experiment~A, macro-F1 is numerically identical to weighted F1
because all five classes have equal support. Per-class inspection
nonetheless reveals that even the best model is uneven across labels:
\texttt{gpt-oss:20b} reaches recall 0.885 on
\texttt{credit\_reporting} but only 0.625 on
\texttt{debt\_collection}. This motivates the macro-F1 reporting in
Experiment~B, where minority-class behaviour can otherwise be hidden
by the majority class. Detailed per-class precision, recall, and F1
for the leading models are reported in Table~\ref{tab:per_class}.

Fig.~\ref{fig:f1_bar} visualises the F1-score distribution across
models. Reasoning models (dark blue) cluster at the top while
standard models show greater spread. The confusion matrix for the
best-performing model (\texttt{gpt-oss:20b}) is shown in
Fig.~\ref{fig:conf_best}. The model handles \texttt{credit\_reporting}
and \texttt{retail\_banking} well but shows some confusion between
\texttt{mortgages\_and\_loans} and \texttt{credit\_card}. Fig.~\ref{fig:conf_phi4} shows the confusion matrix for
\texttt{phi4:14b}, the best standard model: its strong diagonal
confirms reliable five-class classification with only one
off-label prediction, and minor confusion between
\texttt{mortgages\_and\_loans} and \texttt{credit\_card}.
The confusion matrix for \texttt{llama3.1:8b}
(Fig.~\ref{fig:conf_worst}) illustrates the impact of heavy
off-label prediction: a large fraction of samples are assigned
invented categories, degrading performance across all five classes.

\begin{figure}[tb]
  \centering
  \includegraphics[width=\columnwidth]{figures/f1_bar.png}
  \caption{Weighted F1-scores for all models in Experiment~A (Balanced).
           Reasoning models are shown in dark blue; standard models in
           light blue.}
  \label{fig:f1_bar}
\end{figure}

\begin{figure}[tb]
  \centering
  \includegraphics[width=0.92\columnwidth]{figures/confusion_gpt-oss_20b.png}
  \caption{Confusion matrix for \texttt{gpt-oss:20b} (best model,
           Experiment~A). Rows are true labels; columns are
           predicted labels.}
  \label{fig:conf_best}
\end{figure}

\begin{figure}[tb]
  \centering
  \includegraphics[width=0.92\columnwidth]{figures/confusion_phi4_14b.png}
  \caption{Confusion matrix for \texttt{phi4:14b} (best standard model,
           Experiment~A). Strong diagonal with minor confusion between
           \texttt{mortgages\_and\_loans} and \texttt{credit\_card},
           and only one off-label prediction.}
  \label{fig:conf_phi4}
\end{figure}

\begin{figure}[tb]
  \centering
  \includegraphics[width=0.92\columnwidth]{figures/confusion_llama3.1_8b.png}
  \caption{Confusion matrix for \texttt{llama3.1:8b} (worst model,
           Experiment~A). The high off-label rate (57 predictions)
           manifests as diffuse error across all classes.}
  \label{fig:conf_worst}
\end{figure}

\subsection{Experiment B: Imbalanced Evaluation}

Table~\ref{tab:results_imbalanced} presents results for Experiment~B.
All eight models show numerically higher accuracy and F1 compared to
Experiment~A. \texttt{gpt-oss:20b} again leads with accuracy~0.8040
and F1~0.8042. Notably, the standard model \texttt{phi4:14b} reaches
a near-identical score (accuracy~0.7980, F1~0.7987), while
\texttt{qwen3:8b} also matches it (accuracy~0.7980, F1~0.7950).
The table therefore shows that the imbalanced setting changes the
absolute score level more than the identity of the leading models:
\texttt{gpt-oss:20b}, \texttt{phi4:14b}, \texttt{qwen3:8b}, and
\texttt{deepseek-r1:14b} remain the strongest group. However, the
large gains for \texttt{llama3.1:8b} and \texttt{mistral:7b} should
not be interpreted as genuine robustness improvements, because both
models still produce many off-label outputs and benefit from the
dominant \texttt{credit\_reporting} class.

\begin{table}[!t]
  \renewcommand{\arraystretch}{1.3}
  \caption{Experiment B Results — Imbalanced ($N = 1{,}000$)}
  \label{tab:results_imbalanced}
  \centering
  \scriptsize
  \setlength{\tabcolsep}{2pt}
  \resizebox{\columnwidth}{!}{%
  \begin{tabular}{@{}llrrrrrrc@{}}
    \toprule
    \textbf{Model} & \textbf{Type} & \textbf{Acc.} & \textbf{Prec.} & \textbf{Rec.} & \textbf{F1} & \textbf{Macro-F1} & \textbf{Lat. (s)} & \textbf{OL} \\
    \midrule
    \texttt{gpt-oss:20b}      & R & 0.804 & 0.806 & 0.804 & 0.804 & 0.724 & 2.39 & 0 \\
    \texttt{phi4:14b}         & S & 0.798 & 0.803 & 0.798 & 0.799 & 0.700 & 0.54 & 1 \\
    \texttt{qwen3:8b}         & R & 0.798 & 0.799 & 0.798 & 0.795 & 0.703 & 3.60 & 0 \\
    \texttt{deepseek-r1:14b}  & R & 0.790 & 0.802 & 0.790 & 0.794 & 0.702 & 5.72 & 0 \\
    \texttt{deepseek-r1:8b}   & R & 0.771 & 0.777 & 0.771 & 0.768 & 0.663 & 4.41 & 4 \\
    \texttt{gemma2:9b}        & S & 0.762 & 0.785 & 0.762 & 0.762 & 0.661 & 0.66 & 0 \\
    \texttt{llama3.1:8b}      & S & 0.687 & 0.782 & 0.687 & 0.709 & 0.556 & 0.56 & 49 \\
    \texttt{mistral:7b}       & S & 0.686 & 0.777 & 0.686 & 0.703 & 0.536 & 0.23 & 36 \\
    \bottomrule
    \multicolumn{9}{l}{\footnotesize R=Reasoning, S=Standard, OL=Off-label, Lat.=Mean latency}
  \end{tabular}
  }
\end{table}

The relative model ranking is broadly preserved compared to
Experiment~A. However, as discussed in Section~\ref{sec:discussion},
the improvement in accuracy scores must be interpreted carefully in
light of the dominant \texttt{credit\_reporting} class.
The added macro-F1 column makes this clearer: \texttt{gpt-oss:20b}
has weighted F1~0.8042 but macro-F1~0.7237, while
\texttt{llama3.1:8b} has weighted F1~0.7092 but macro-F1~0.5555.
The gap shows that weighted metrics substantially overstate
minority-class performance in the imbalanced setting. The
F1-score distribution across all Experiment~B models mirrors
Experiment~A in its ranking (Fig.~\ref{fig:comparison_ab} right bars),
with \texttt{gpt-oss:20b}, \texttt{phi4:14b}, and \texttt{qwen3:8b}
again forming the leading group.

\subsection{Cross-Experiment Comparison}

Fig.~\ref{fig:comparison_ab} provides a side-by-side view of accuracy
across both experiments. All models improve numerically from
Experiment~A to Experiment~B; the shift is most pronounced for
\texttt{llama3.1:8b} ($\Delta$acc~$= +0.100$)
and \texttt{mistral:7b} ($\Delta$acc~$= +0.090$), the two models with the
highest off-label counts. In terms of weighted F1 the corresponding
gains are even larger ($+0.113$ and $+0.119$ respectively). This
pattern suggests that under imbalanced conditions, a model defaulting
to the majority class can still achieve a deceptively high aggregate
score.

\begin{figure}[tb]
  \centering
  \includegraphics[width=\columnwidth]{figures/comparison_A_vs_B.png}
  \caption{Side-by-side comparison of accuracy across
           Experiment~A (Balanced, blue) and Experiment~B (Imbalanced, red)
           for all models. The dashed reference line is labelled
           o1 (0.777), the best paid-API result from the reference paper. The
           upward shift in Experiment~B partially reflects
           class imbalance rather than genuine model improvement.}
  \label{fig:comparison_ab}
\end{figure}

\subsection{Comparison with Paid-API Models}

Fig.~\ref{fig:vs_reference} compares all eight of our open-source models
with the five best reported results from Roumeliotis~et~al.\
\cite{roumeliotis2025think}. Table~\ref{tab:vs_reference} summarises
the key comparison points. Our top model, \texttt{gpt-oss:20b},
achieves a higher F1 (0.7451) than the GPT-4o score reported in the
reference benchmark (0.7358) while requiring no API expenditure.
\texttt{phi4:14b} and \texttt{qwen3:8b} both approach the GPT-4o
performance tier. Because the reference paper does not publish the
exact sampled complaint identifiers, this comparison should be read
as protocol-matched rather than paired on identical instances. The
gap to o1 (F1~0.7764) and GPT-4.5 (F1~0.7763) remains non-trivial,
suggesting that the reasoning depth of frontier proprietary models
still provides an edge for this structured task. Nevertheless, the
cost differential is stark: the reference paper reported o1 costs of
\$24.24 per thousand predictions, versus \$0.00 for all our models.
Thus, Table~\ref{tab:vs_reference} supports a more nuanced claim than
simple ``open-source wins'' or ``paid APIs win'': the best local model
does not reach the strongest paid reasoning models, but it does exceed
the reported GPT-4o baseline and removes recurring API cost entirely.
This is the central deployment trade-off of the paper. For low-budget,
privacy-sensitive, or offline environments, \texttt{gpt-oss:20b} offers
a competitive zero-shot option; for settings where the last few F1
points matter more than cost or data locality, frontier paid reasoning
models still retain an advantage.

\begin{table}[!t]
  \renewcommand{\arraystretch}{1.3}
  \caption{Comparison with Reference Paper (Balanced, $N = 1{,}000$).
           Paid-API results sourced from Roumeliotis~et~al.\ \cite{roumeliotis2025think}.}
  \label{tab:vs_reference}
  \centering
  \scriptsize
  \setlength{\tabcolsep}{2pt}
  \resizebox{\columnwidth}{!}{%
  \begin{tabular}{@{}llrrl@{}}
    \toprule
    \textbf{Model} & \textbf{Source} & \textbf{Accuracy} & \textbf{F1} & \textbf{Cost/1k} \\
    \midrule
    o1                          & Paid API & 0.777 & 0.7764 & \$24.24 \\
    GPT-4.5                     & Paid API & 0.774 & 0.7763 & \$12.32 \\
    Claude~3.7~Sonnet           & Paid API & 0.765 & 0.7626 & \$0.64 \\
    Gemini~2.0~Flash            & Paid API & 0.758 & 0.7586 & \$0.04 \\
    GPT-4o                      & Paid API & 0.736 & 0.7358 & \$0.45 \\
    \midrule
    \texttt{gpt-oss:20b}        & Local (ours)                   & 0.745 & 0.7451 & \$0.00 \\
    \texttt{deepseek-r1:14b}    & Local (ours)                   & 0.716 & 0.7162 & \$0.00 \\
    \texttt{phi4:14b}           & Local (ours)                   & 0.709 & 0.7095 & \$0.00 \\
    \texttt{qwen3:8b}           & Local (ours)                   & 0.707 & 0.7066 & \$0.00 \\
    \bottomrule
  \end{tabular}
  }
\end{table}

\begin{figure}[tb]
  \centering
  \includegraphics[width=\columnwidth]{figures/vs_reference.png}
  \caption{Comparison of all eight open-source models (Experiment~A)
           against five paid-API baselines from Roumeliotis~et~al.\
           \cite{roumeliotis2025think}. Horizontal bars show accuracy.}
  \label{fig:vs_reference}
\end{figure}

\subsection{Latency Analysis}

Fig.~\ref{fig:latency} plots accuracy against mean prediction latency
for all models in Experiment~A. Standard models are notably faster:
\texttt{mistral:7b} completes a prediction in 0.24~s on average,
while \texttt{phi4:14b} requires only 0.55~s despite being twice the
parameter count. Reasoning models are slower by one to two orders of
magnitude relative to the smallest standard models:
\texttt{deepseek-r1:14b} takes 5.92~s per prediction due to its
chain-of-thought trace. The latency–accuracy trade-off is therefore
sharply model-type dependent: \texttt{phi4:14b} achieves the best
accuracy-to-latency ratio among all models, requiring just 0.55~s
per prediction while reaching F1~0.7095.

\begin{figure}[tb]
  \centering
  \includegraphics[width=\columnwidth]{figures/latency_vs_accuracy.png}
  \caption{Accuracy vs.\ mean prediction latency (seconds) for all
           models in Experiment~A. Reasoning models (diamonds) tend
           to be slower but more accurate; standard models (circles)
           vary widely. \texttt{phi4:14b} achieves the best
           accuracy-to-latency ratio.}
  \label{fig:latency}
\end{figure}

\subsection{Per-Class Breakdown for Leading Models}
\label{sec:per_class_metrics}

Table~\ref{tab:per_class} reports per-class F1 for the three leading
models in each experiment, confirming that aggregate weighted F1 masks
important minority-class variation. In Experiment~B all three models
jump to ${\geq}0.87$ F1 on \texttt{credit\_reporting} while minority
classes remain in the 0.57–0.76 range, confirming the majority-class
bias effect discussed in Section~\ref{sec:discussion}.

\begin{table}[tb]
  \renewcommand{\arraystretch}{1.2}
  \caption{Per-Class F1 for Leading Models (Exp.~A and~B)}
  \label{tab:per_class}
  \centering
  \scriptsize
  \setlength{\tabcolsep}{3.5pt}
  \begin{tabular}{@{}llccccc@{}}
    \toprule
    \textbf{E} & \textbf{Model} &
    \textbf{Retail} &
    \textbf{CreditRep.} &
    \textbf{Card} &
    \textbf{Mortg.} &
    \textbf{Debt} \\
    \midrule
    A & \texttt{gpt-oss:20b}     & 0.8098 & 0.7195 & 0.7154 & 0.7902 & 0.6906 \\
    A & \texttt{phi4:14b}        & 0.7259 & 0.7054 & 0.6998 & 0.7457 & 0.6705 \\
    A & \texttt{qwen3:8b}        & 0.7700 & 0.6870 & 0.6613 & 0.7740 & 0.6409 \\
    \midrule
    B & \texttt{gpt-oss:20b}     & 0.7771 & 0.8810 & 0.6316 & 0.7179 & 0.6109 \\
    B & \texttt{phi4:14b}        & 0.7006 & 0.8831 & 0.6182 & 0.6624 & 0.6378 \\
    B & \texttt{qwen3:8b}        & 0.7595 & 0.8801 & 0.5750 & 0.7006 & 0.5975 \\
    \bottomrule
    \multicolumn{7}{l}{\footnotesize E = Experiment; Mortg. = Mortgages \& Loans.}
  \end{tabular}
\end{table}

% =============================================================
\section{Discussion}
\label{sec:discussion}

\subsection{Reasoning vs.\ Standard Models}

The most consistent finding across both experiments is that reasoning
models dominate the upper performance tier.
Three of the four reasoning models appear in the top four positions of
Experiment~A (\texttt{gpt-oss:20b}, \texttt{deepseek-r1:14b},
\texttt{qwen3:8b}), with the standard \texttt{phi4:14b} placed third. This pattern aligns with the findings of
Roumeliotis~et~al.\ \cite{roumeliotis2025think}, who similarly found
that o1 and o3-mini outperformed GPT-4o and other non-reasoning
counterparts, and with the broader evidence that chain-of-thought
prompting elicits more reliable structured outputs from LLMs
\cite{wei2022chainofthought,guo2025deepseekr1}. The underlying
mechanism is plausibly that reasoning allows the model to decompose
ambiguous narratives before committing to a label, reducing
noise-driven misclassifications.

However, the advantage is not universal. \texttt{deepseek-r1:8b}
(Reasoning, 8B parameters, F1~0.6736) is outperformed by
\texttt{phi4:14b} (Standard, 14B, F1~0.7095) and is comparable to
\texttt{gemma2:9b} (Standard, 9B, F1~0.6728). This suggests that
\emph{model scale} interacts strongly with the reasoning benefit:
at 8B parameters, reasoning distillation may not provide enough
representational capacity to overcome the advantage of a larger,
high-quality standard model.

\subsection{Off-Label Prediction Behaviour}

Off-label predictions — responses where the model returned a label
not in the prescribed five categories — are a direct measure of
instruction-following reliability. The two models with the worst
F1 scores (\texttt{llama3.1:8b}:~57; \texttt{mistral:7b}:~54)
also exhibit the highest off-label counts in Experiment~A. This is
not coincidental: off-label responses are counted as prediction
errors, directly reducing accuracy and recall. In Experiment~A, all four
reasoning models produce at most one off-label prediction each, and
\texttt{gpt-oss:20b}, \texttt{deepseek-r1:14b}, and
\texttt{qwen3:8b} produce zero.
The better instruction-following of reasoning models reflects the
structured nature of chain-of-thought: the model explicitly reasons
about which label is appropriate before formatting its output,
reducing the likelihood of hallucinating a novel category.

The raw status logs show that off-label behaviour has two distinct
forms. \texttt{llama3.1:8b} primarily fails at the formatting/parsing
level, producing 57 parse errors in Experiment~A and 49 in
Experiment~B. In contrast, \texttt{mistral:7b} usually returns
parseable JSON but invents labels outside the allowed taxonomy,
producing 54 valid-but-out-of-vocabulary labels in Experiment~A and
36 in Experiment~B. Smaller counts for \texttt{phi4:14b},
\texttt{gemma2:9b}, and \texttt{deepseek-r1:8b} are also
valid-but-out-of-vocabulary labels such as \texttt{unknown} or
\texttt{health\_insurance}. This distinction matters operationally:
parse errors indicate output-format failure, whereas out-of-vocabulary
labels indicate failure to stay within the regulatory taxonomy.

Notably, the off-label counts for \texttt{mistral:7b} and
\texttt{llama3.1:8b} improve substantially in Experiment~B
(36 and 49 respectively, compared to 54 and 57 in Experiment~A).
This may reflect a simpler classification decision under
imbalanced conditions, where the dominant \texttt{credit\_reporting}
class provides an easier default that is more likely to be
formatted correctly.

\subsection{Latency–Accuracy Trade-off}

From a practical deployment perspective, the latency results present
an important consideration. The fastest model, \texttt{mistral:7b}
(0.24~s), achieves only F1~0.5842 — the worst result in Experiment~A.
\texttt{phi4:14b}, at 0.55~s per prediction and F1~0.710, offers
the most favourable accuracy-to-latency ratio of any model tested.
For applications requiring both real-time response and reasonable
accuracy, \texttt{phi4:14b} is the recommended choice.

Reasoning models are consistently slower: \texttt{deepseek-r1:14b}
takes 5.92~s per prediction. At 1,000 predictions, this translates
to roughly 1.6~hours of total inference time on a single V100. For
large-scale regulatory applications processing millions of complaints,
the time budget would need to be considered alongside model quality.
All reported latency figures for reasoning models include the full
chain-of-thought generation trace; the final JSON label is extracted
from the end of this trace via the regex described in
Section~\ref{sec:setup}. Latency comparisons between reasoning and
standard models should therefore be interpreted as end-to-end
wall-clock cost rather than pure decoding speed.

\subsection{Interpreting the Imbalanced Results}

The across-the-board increase in accuracy and F1 from Experiment~A
to Experiment~B demands careful interpretation. While a genuine
improvement is possible — models may be aided by the greater
frequency of \texttt{credit\_reporting} training signal in
their pre-training corpora, which coincides with its frequency in
the CFPB data — a more parsimonious explanation exists:
\emph{class imbalance inflates aggregate accuracy metrics}.

When 63.6\,\% of the test set belongs to a single class, a model
that classifies most ambiguous narratives as \texttt{credit\_reporting}
will achieve a high weighted accuracy even if its per-class performance
is poor. The particularly pronounced gains for \texttt{llama3.1:8b}
(+11.3 F1 points) and \texttt{mistral:7b} (+11.9 F1 points) — the
two highest off-label producers in Experiment~A — are consistent
with this mechanism: under imbalanced conditions, a confused model
benefits disproportionately from majority-class guesses. We therefore
caution against using Experiment~B accuracy as a simple indicator of
model quality and recommend that future evaluations report per-class
precision and recall alongside aggregate metrics.

\subsection{Comparison with the Reference Paper}

Our results confirm that competitive zero-shot classification is
achievable with open-source, locally deployed models. The
\texttt{gpt-oss:20b} model exceeds the reported GPT-4o benchmark on
both accuracy and F1 under the matched five-category protocol, while
incurring no API fee. This is a meaningful result: GPT-4o costs
\$0.45 per thousand predictions in the reference benchmark, meaning
that at moderate scale (one million complaints per year), local
deployment would avoid roughly \$450 in API fees while keeping
classification quality in the same performance tier. This should not
be read as computation being free: local inference still consumes GPU
time, electricity, system administration effort, and institutional HPC
capacity. The practical claim is therefore about eliminating recurring
third-party API charges and data transfer, not about eliminating all
operational cost.

The remaining gap to the frontier paid models (o1, GPT-4.5,
Claude~3.7) is approximately 3–4 F1 points. Whether this gap
is consequential depends on the application. For initial triage
and routing, an F1 around 0.74–0.75 may be fully adequate.
For compliance-critical decisions, the gap warrants attention —
though it could plausibly be closed by fine-tuning
\cite{gullu2025llama,li2024customerservice} or ensemble methods.

\subsection{Limitations}

This study has several limitations. First, the balanced comparison
matches the reference protocol but not necessarily the exact sampled
instances, since the original sample identifiers are unavailable.
Second, the 500-character length filter reduces inference cost but
also shifts the class distribution, especially increasing the share of
\texttt{credit\_reporting}; future work should test whether the same
patterns hold for longer narratives. Third, lowercasing and special-
character removal may discard useful cues such as account-number
patterns, capitalization, dates, or punctuation; the effect of less
aggressive LLM-oriented preprocessing remains untested. Fourth, all
experiments use
$N{=}1{,}000$ examples, one random seed, and a five-class taxonomy;
larger temporal slices, multiple resampled test sets, and richer
taxonomies may reveal additional drift and rarer category confusions.
Fifth, all local models are evaluated with a single zero-shot prompt.
Prompt wording, output-format instructions, few-shot examples, or
self-consistency decoding may change the relative ranking of models,
particularly for smaller standard models with weaker instruction
following. Sixth, all local models are evaluated in their q4 Ollama
quantised form, so the reported scores conflate model family, model
size, and quantisation effects. Seventh, we cannot rule out pre-training
contamination: open-weight models may have encountered CFPB complaints
or derivative datasets during pre-training, although no task-specific
fine-tuning is used in our experiments. Eighth, the reported \$0 cost
refers to API expenditure, not to the full economic cost of GPU
hardware, electricity, or HPC operation. Finally, the main tables
report weighted aggregate metrics; for operational deployment,
macro-F1 and per-class precision/recall should be inspected to avoid
masking poor minority-class performance.

\subsection{Infrastructure Contribution}

A secondary contribution of this work is the documented HPC
deployment pipeline. The combination of SLURM job scheduling,
offline Ollama serving, and a standard-library-only Python inference
script provides a template that can be applied to other NLP
classification tasks on shared research clusters without requiring
pip-installed dependencies on the compute node. This reproducibility
infrastructure lowers the barrier for other researchers to benchmark
local LLMs on domain-specific tasks.

% =============================================================
\section{Conclusion}
\label{sec:conclusion}

We have presented a comprehensive zero-shot benchmark of eight
open-source large language models on the CFPB consumer financial
complaint classification task, conducted entirely at zero API cost
on the TRUBA ARF HPC cluster. Our results demonstrate that the
best open-source model, \texttt{gpt-oss:20b}, achieves a weighted
F1 of 0.7451 on a balanced evaluation, exceeding the reported paid
GPT-4o score (F1~0.7358) from the reference paper by
Roumeliotis~et~al.\ \cite{roumeliotis2025think} under a matched
protocol. Reasoning models generally occupy the strongest performance
tier, although the large standard model \texttt{phi4:14b} (F1~0.7095)
offers the best latency–accuracy trade-off among all models tested.

Our secondary imbalanced experiment reveals important evaluation
dynamics: aggregate accuracy improvements under imbalanced
conditions are partially attributable to majority-class bias
rather than genuine model improvement, underscoring the need for
per-class reporting in real-world settings.

From a practical standpoint, these findings validate the viability
of self-hosted LLM deployments for financial NLP tasks: institutions
that cannot share sensitive complaint text with third-party API
providers, or that operate under budget constraints, can now achieve
competitive classification quality using open-weight models on
standard GPU hardware.

Future work will explore five directions.
First, a \emph{quantisation sensitivity study} will evaluate all eight
models at q4, q8, and fp16 precision to determine whether the accuracy
gains from higher numerical precision justify the increased VRAM cost — a
sub-study not attempted in the reference paper.
Second, a \emph{prompt and sampling sensitivity study} will repeat the
balanced and imbalanced evaluations across multiple random seeds and
prompt variants, reporting confidence intervals and paired significance
tests beyond the post-hoc McNemar analysis presented here.
Third, and most directly motivated by the present results, we will
\emph{fine-tune the same eight open-weight models} on the CFPB
training split and compare their supervised performance against their
zero-shot counterparts reported here. This paired comparison will
quantify the exact cost of removing the annotation requirement and
clarify whether the 3–4 point F1 gap to frontier paid APIs can be
closed by task-specific supervised training.
Fourth, \emph{long-narrative robustness} will be examined by relaxing
the 500-character filter and measuring whether additional complaint
context improves classification or instead increases prompt length and
latency costs.
Fifth, the benchmark will be \emph{extended to the full multi-year
CFPB corpus} with a richer category taxonomy to assess scalability
and temporal robustness of the open-source inference pipeline.

% ---- Acknowledgements ----------------------------------------
\section*{Acknowledgment}
The numerical calculations reported in this paper were fully performed
at T{\"U}B\.{I}TAK ULAKB\.{I}M, High Performance and Grid Computing
Center (TRUBA resources). The authors gratefully acknowledge
T{\"U}B\.{I}TAK ULAKB\.{I}M for providing access to the TRUBA ARF
cluster and the \texttt{akya-cuda} GPU partition used in this study.

% ---- Conflict of Interest ------------------------------------
\section*{Conflict of Interest}
The authors declare that they have no conflict of interest.

% ---- Funding Statement ---------------------------------------
\section*{Funding}
This research received no external funding.

% ---- Code Availability ---------------------------------------
\section*{Code Availability}
All source code, preprocessing scripts, and SLURM job files used in
this study are publicly available at:
\url{https://github.com/muhammettan28/ConsumerComplaint}

% ---- References ---------------------------------------------
% Balance final two columns at approximately ref #20
\IEEEtriggeratref{16}
\bibliographystyle{IEEEtran}
\bibliography{references}

\end{document}
